\documentclass[11pt]{article}
\usepackage{amsmath,amssymb}
\title{Room 10 (P4 D-algebraicity): moderator summary}
\date{2026-09-26}
\begin{document}
\maketitle

One seat (Hilbert) attempted to construct the missing disc-analog Nevanlinna growth theorem
(ingredient (b) of \texttt{P4\_note.tex}) from first principles, since neither of the two prior
rooms on P4 (Room 4, and Track A4 before it) had attempted building this machinery directly.

\textbf{Hilbert's derivation.} Re-examining the classical proof of Nevanlinna's
logarithmic-derivative lemma (Poisson--Jensen, two radii $r<R$), the proof itself is $r/R$-local
and never actually uses $R\to\infty$ -- the plane-vs-disc distinction only enters afterward, in how
$R(r)$ is chosen and the growth scale used. Transplanting mechanically ($r\to\infty \leftrightarrow
r\to1^-$) gives a candidate disc log-derivative lemma and, combined with the fact that polynomial
ODE coefficients have $T(r,\text{coeff})=O(1)$ on a compact disc (rather than growing like $\log r$
on the plane), yields a candidate bound $T(r,y)=O(\log(1/(1-r)))$ for any $y$ meromorphic on the
disc solving a fixed algebraic ODE of order $k$, degree $d\ge2$. This is a real, if unreviewed,
piece of original derivation.

\textbf{The bookkeeping question, resolved.} The comparison against $U$'s already-established
$n(r,U)=\Theta(1/(1-r))$ pole count (Room 4) hinges entirely on which counting-function convention
is standard in disc Nevanlinna theory: $N(r,f)=\int_0^r n(t)\,dt/t$ (the plane integral, restricted
to $r<1$) or $N(r,f)=\int_0^r n(t)\,dt/(1-t)$ (a disc-native weighting that would make the boundary
play the role $\infty$ plays in the plane theory). Hilbert flagged this as unresolved and
UNVERIFIED-CITATION, without library access. \textbf{Checked here against primary sources (web
search, cross-referencing the standard definition and Tsuji's disc extension): the standard
convention is $N(r,a)=\int_0^r[n(t,a)-n(0,a)]/t\,dt+n(0,a)\log r$ -- the SAME $dt/t$ weighting as
the plane theory, not a $dt/(1-t)$ disc-native one.} This is Hilbert's ``alternative convention,''
which the seat's own analysis already correctly identified as giving an \textbf{inconclusive}
comparison (same order, no violation) rather than the ``disc-native'' convention that would have
closed P4.

\textbf{Room verdict: NO PROGRESS, but a real sub-question genuinely resolved.} The candidate disc
growth theorem, even if it survives full scrutiny, does \emph{not} close P4 under the standard
convention -- the comparison is inconclusive, exactly as one of Hilbert's two branches anticipated.
This settles a real ambiguity (rather than leaving it open) and rules out a specific hoped-for path
to closure. P4 stands exactly where \texttt{P4\_note.tex}/Room 4 left it: OPEN, with the correct
theorem family named, ingredient (a) now precisely known ($O(1/(1-r))$, not superpolynomial), and
ingredient (b)'s standard-convention version now known not to be sufficient even if constructed.
The candidate log-derivative-lemma transplant itself (genuinely new, unreviewed) is left in
\texttt{room10\_seat\_hilbert.tex} for anyone who wants to pursue a non-standard/alternative growth
measure in the future, though that would need independent justification for why it's the right
notion of ``growth'' to use.

\end{document}
